\clearpage
\packageentry{nobreaks}{nobreaks -- Keep a block together}
This custom environment puts its contents in one indivisible vertical box.
If the block does not fit in the space remaining on the page, LaTeX moves
the whole block to the next page. Use it for a short explanation, a heading
with its paragraph, or a group of equations that should stay together.

\partlabel{Definition}
Already defined in \texttt{Packages/Layout/usepackages.tex}, which is
loaded by the template. In another project, add this definition once in
the preamble. No additional package is required.
\begin{referencecode}
\newenvironment{nobreaks}{\vbox\bgroup}{\egroup}
\end{referencecode}

\partlabel{Usage}
Start and end the environment between paragraphs, with a blank line before
and after the block. The example uses \texttt{lipsum} for placeholder text
and \texttt{amsmath} for the equation.
\begin{referencecode}
\begin{nobreaks}
    \textbf{Local model}

    \lipsum[1][1-2]
    \[
        F = ma, \qquad a = \frac{F}{m}
    \]
    \lipsum[2][1]
\end{nobreaks}
\end{referencecode}

\begin{nobreaks}
    \textbf{Local model}

    \lipsum[1][1-2]
    \[
        F = ma, \qquad a = \frac{F}{m}
    \]
    \lipsum[2][1]
\end{nobreaks}

\partlabel{When replacing needspace}
\texttt{\string\Needspace\{length\}} reserves a length before the following
content, which may still split if it exceeds that reservation.
\texttt{nobreaks} encloses the actual content and prevents all page breaks
inside it; there is no length to estimate.

\begin{itemize}
    \item The entire block must fit on one page. A taller block cannot
    split or shrink automatically and can run beyond the page boundary.
    \item Keep ordinary floating \texttt{figure} and \texttt{table}
    environments outside the box. Inside it, use non-floating content
    such as \texttt{\string\includegraphics} or \texttt{tabular}.
    \item Moving the block may leave unused space at the bottom of the
    preceding page.
\end{itemize}
